% @Claude and @codex NEVER edit this without explicit instruction to do so. If adding a figure, table or other non-text content, DO NOT edit the text without instruction to do so, even to add a reference to the new content.
% One sentence per line!
% @claude: make edits in red
% @codex: make edits in green
\documentclass{article} % For LaTeX2e
\usepackage{iclr2026_conference,times}

\usepackage{hyperref}
\usepackage{url}
\usepackage{xcolor}
\usepackage{graphicx}
\usepackage{subcaption}
\usepackage{placeins}
\usepackage{booktabs}
\usepackage{amsmath,amssymb}
\usepackage{dsfont} % \mathds{1} indicator (Type 1 fonts, unlike bbm)
\usepackage{textcomp}
\usepackage{listings}
\lstdefinestyle{promptstyle}{
  basicstyle=\ttfamily\scriptsize,
  breaklines=true,
  breakautoindent=false,
  columns=fullflexible,
  frame=single,
  framesep=4pt,
  xleftmargin=0pt,
  showstringspaces=false,
  keepspaces=true,
}
\usepackage{amsthm}
\newtheorem{proposition}{Proposition}
\theoremstyle{plain}
\newtheorem{theorem}{Theorem}
\newtheorem{lemma}[theorem]{Lemma}
\newtheorem{corollary}[theorem]{Corollary}
\theoremstyle{definition}
\newtheorem{definition}[theorem]{Definition}
\newtheorem{remark}[theorem]{Remark}

%% Commands for TeXCount
%TC:macro \cite [option:text,text]
%TC:macro \citep [option:text,text]
%TC:macro \citet [option:text,text]
%TC:envir table 0 1
%TC:envir tabular [ignore] word
%TC:envir displaymath 0 word
%TC:envir math 0 word
%TC:envir comment 0 0

%% Red text for drafting notes/highlights: \red{...}
\newcommand{\red}[1]{\textcolor{red}{#1}}
\newcommand{\green}[1]{\textcolor{green}{#1}}
\definecolor{darkgreen}{RGB}{0,100,0}
\newcommand{\darkgreen}[1]{\textcolor{darkgreen}{#1}}

%% Dark blue for the Notes section
\definecolor{darkblue}{RGB}{0,0,139}

%% Indicator function: \ind{d \not\models c_i}
\newcommand{\ind}[1]{\mathds{1}\!\left[#1\right]}

%% Inner product: \inner{q}{d}
\newcommand{\inner}[2]{\left\langle #1, #2 \right\rangle}

\title{Partial Credit Improves Retrieval Embeddings on Conjunctive Queries}

% Authors must not appear in the submitted version. They should be hidden
% as long as the \iclrfinalcopy macro remains commented out below.
\author{Anonymous Authors \\
Institution \\
\texttt{email@example.com}
}

%\iclrfinalcopy % Uncomment for camera-ready version, but NOT for submission.
\begin{document}

\maketitle

\begin{abstract}
% \red{TODO: abstract.}

As users refine search queries, especially in e-commerce, they tend to add requirements as conjunctions, making each query more specific than the last.
Prior work has shown that embedding-based retrieval accuracy degrades as the number of conditions grows, motivating methods that remain robust under conjunctive queries.
One limitation in existing training paradigms is the binary training objective; hard negatives that subtly conflict with the query are considered just as incorrect as easy, obviously irrelevant negatives.
To leverage the wasted information, we present the Graded Objective for Learning Distinctions, or \textsc{gold} mining, a hard negative mining strategy that generates hard negatives with known violations of query conditions.
Additionally, we contribute \textsc{gold} reweighting, a loss function modification strategy that assigns partial credit to hard negatives that nearly satisfy the query.
We demonstrate that both \textsc{gold} mining and \textsc{gold} reweighting are compatible with four common existing loss function families.
Applying both outperform state-of-the-art hard negatives mining in seven evaluation settings across four loss functions and two modalities.
Our theoretical analysis establishes that under suitable query and comparison coverage assumptions, uniformly low loss on the \textsc{gold} Margin-MSE objective implies approximate affine recoverability of attributes, a property not guaranteed satisfiable by the standard objective.


% Given limited assumptions, theoretical analysis of \green{Ours} applied to Margin-MSE shows that the graded loss objective can be satisfied when the query vector serves as an affine probe for the conjunctive conditions in the query, and that .

% Under explicit norm and scale assumptions, the hard-negative Margin-MSE objective admits low-loss cosine query representations formed from signed sums of the vector parts of affine attribute probes.
% Under a coverage assumption, low count loss on a query pair differing in one condition implies approximate affine recovery of the deleted attribute, whereas zero binary Margin-MSE loss on the same query family need not.



\end{abstract}

\section{Introduction}

Embedding-based retrieval aims to recall a match set of semantically relevant documents satisfying the user's query.
Such embeddings are often used in retrieval tasks, such as retrieval augmented generation (RAG), product search, and image search, where fine-grained matching of details between the query and the document determines relevance.
As users' refine their queries, they tend to increase in specificity over time by filtering out more candidate results.
Such user behavior is demonstrated in employee recruitment  \citep{kaya2023understanding}, web search \citep{jansen2007patterns}, and product search \citep{papenmeier2020modern}.
These refined queries are often conjunctive: a result must satisfy every condition, including negated conditions \citep{merra2023improving}



However, bi-encoders are known to struggle with exactly the operations that such queries require.
These include negation \citep{zhang2025excluir, weller2024nevir}, set operations \citep{malaviya-etal-2023-quest}, simultaneous constraint satisfaction \citep{lu2025multiconir, killingback2025benchmarking}, and attribute binding \citep{yuksekgonul2022and}.
We hypothesize that user query behavior may suffer from a chicken-and-egg problem; users do not use complex queries because they expect models to fail, while practitioners underprioritize complex queries because they do not appear in user data.


Although Amazon ESCI \cite{reddy2022shopping} includes simple multi-attribute queries, it offers limited coverage of complex conjunctions, motivating our dataset’s systematic construction of queries combining multiple required and forbidden attributes.
This neglects an important dynamic in the user experience of search systems wherein queries tend to gain attributes over time as users refine them. 
\citet{hirsch2020query} find that over 50\% of queries on eBay occur during a reformulation session, with queries growing in length over time.
Longer, rewritten queries were associated with increased purchase rates, indicating that support for highly specific queries should be a priority.

Thus, we identify dense retrieval for conjunctive queries as understudied across three axes: data, methods, and evaluations.
Specifically, we concentrate on improving conjunctive queries --- those seeking documents that \textbf{have} attributes a, b, c ... and \textbf{do not have} attributes x, y, z, ... . 
This focus is further grounded in patterns in user queries and the design of pre-existing non-embedding search systems.
For example, the conjunctive filtering paradigm already exists in product search in the form of attribute filters on e-commerce sites, such as ``screen size $\geq 38$ cm'' for laptops.
% Users find such categories useful both in searching for and understanding products. 
Unfortunately, attribute ontologies are usually fixed, failing when a user desires a product with or without a novel attribute.
% Because unique and novel attributes differentiate products, sellers are well motivated to include them in product descriptions to differentiate their products. 
As a result, natural language search systems, such as dense retrieval, are better suited for the long tail of highly specific queries.
Similar use cases exist in other retrieval tasks, including for images, videos, resumes, patents, and real estate.



To address this gap, we introduce \textbf{G}raded \textbf{O}bjective for \textbf{L}earning \textbf{D}istinctions, or \textsc{gold} mining, a strategy for constructing conjunctive queries for existing document pairs, automatically labeling how many query conditions each hard negative violates.
We combine it with \textsc{gold} reweighting, a broadly applicable loss function modification that assigns partial credit to hard negatives that partially satisfy those conditions.
\textsc{gold} mining uses a generative LLM to extract attributes of each document in structured form, creating a large, semi-synthetic Main dataset.
Then, we construct triplets consisting of a natural language query, positive document, and hard negative document based on the structured attributes.
Leveraging the structured attribute lists allows us to calculate the number of query conditions violated by the hard negative, which we use in the \textsc{gold} reweighting training objective.
% We then use a Hamming distance-based loss function to train the embedding model.
This contrasts with baseline loss functions that treat all negatives as equally wrong.
% For the query ``warm red hoodie'', the document ``warm blue jacket'' is considered just as incorrect as ``gas-powered lawnmower''.
% On the other hand, Ours sets the loss proportional to the squared number of attributes that must be changed to satisfy the query, and much higher if the subject (hoodie vs. lawnmower) must be changed.

% With global e-commerce sales exceeding \$28 trillion USD in 2024 and numerous multimodal reasoning tasks' reliance on fine-grained visual understanding, 

We demonstrate the effectiveness of these strategies on both the text and image modalities.
We choose textual product search due to its economic importance, with annual e-commerce sales exceeding \$28 trillion USD in 2024 \citep{unctadstat2026digital}. 
We choose apparel image search because it requires objective reasoning over fine-grained visual details, a foundational skill in numerous multimodal tasks.
Although we demonstrate the effectiveness of \textsc{gold} mining and reweighting on these two domains, we stress that, in principle, the strategies are highly general.
\textsc{gold} mining can be applied in settings where documents can be reasonably well specified by binary QA (i.e., a list of attributes).
\textsc{gold} reweighting need not rely on weights derived from \textsc{gold} mining specifically, and could instead use, for example, pre-existing structured data.
Our results demonstrate consistent improvements over baselines using \textsc{gold} mining, with further improvement after adding \textsc{gold} reweighting.
% Our results show that \textsc{gold} mining improves recall across 7/8 conditions, especially in the text domain.
% In most cases, this improvement exceeds that of the state-of-the-art NV-Retriever mining algorithm.
% When \textsc{gold} reweighting is applied on top of \textsc{gold} mining, recall consistently improves further.


Our theoretical analysis shows that, under suitable query and comparison coverage assumptions, uniformly low loss on the \textsc{gold} Margin-MSE objective excluding easy negatives implies approximate affine recoverability of attributes, a property not guaranteed even by zero binary-target Margin-MSE loss.
The converse is also true: affine recoverability of attributes permits low loss on the \textsc{gold} Margin-MSE objective under a reasonable scale assumption.
Finally, we show that an embedding of fixed length $d$ can represent exponentially many attributes with low loss.



\section{Related Work}

\paragraph{Dense retrieval and graded supervision.}
Dense bi-encoders enable efficient first-stage retrieval by mapping queries and candidates into a shared embedding space~\cite{reimers-gurevych-2019-sentence,karpukhin-etal-2020-dense}.
% Their effectiveness, however, depends strongly on the selection and supervision of negative examples.
Methods such as ANCE and RocketQA improve training by retrieving challenging negatives, sharing negatives across batches, or denoising mined candidates~\cite{xiong2020approximate,qu-etal-2021-rocketqa}.
Other approaches move beyond uniform positive--negative separation by distilling continuous teacher-score differences, as in Margin-MSE and TAS-Balanced~\cite{hofstatter2020improving,hofstatter2021tas}, or by using relevance- and attribute-dependent margins in metric-learning objectives~\cite{zhou2020ladder,Jeong_2021_ICCV,Zhao2019}.
% \red{todo: discuss ColBERTv2 KL distillation over mined candidates (https://arxiv.org/abs/2112.01488), BGE/C-Pack rank-window ANN mining (https://arxiv.org/abs/2309.07597), NV-Retriever positive-aware negative filtering (https://arxiv.org/abs/2407.15831)}
% These methods demonstrate the value of informative negatives and graded training signals, but generally derive the degree of separation from teacher scores, relevance labels, or broad semantic relationships.
% In contrast, we derive supervision from the discrete structure of the query: candidates are graded by the number of query-relevant feature edits required to make them satisfy its constraints.

\paragraph{Compositional, Boolean, and product retrieval.}
% Product relevance is often compositional, as a relevant item may need to satisfy multiple positive and negative attribute constraints simultaneously.
% Prior product-search methods address this problem through behaviorally informed negative sampling or representations that explicitly model product facets and attribute composition~\cite{Mohan2019,kong2022multiaspect,Choudhary2022}.
Prior product search methods address the problem of compositional requirements through behaviorally informed negative sampling or explicitly representing facets and attribute composition~\cite{Mohan2019,kong2022multiaspect,Choudhary2022}.
A related line of work studies logical composition in dense retrieval.
QUEST considers natural-language queries representing set operations, while BoolQuestions examines conjunction, disjunction, and exclusion through Boolean-aware training and query decomposition~\cite{malaviya-etal-2023-quest,zhang-etal-2024-boolquestions}.
Logical query-embedding methods such as BetaE and ConE instead construct geometric representations that directly support logical operators~\cite{NEURIPS2020_e43739bb,NEURIPS2021_a0160709}.
NevIR shows that bi-encoders are often insensitive to negated meaning, while product-search systems have mitigated related failures through explicit detection and filtering of prohibited attributes~\cite{weller-etal-2024-nevir,Merra2023}.
In the multimodal domain, the CIRR benchmark and methods such as CLIP4CIR study retrieval from queries that combine a reference image with a textual modification~\cite{pmlr-v139-radford21a,Liu_2021_ICCV,Baldrati_2022_CVPR}.
% More recent approaches introduce supervision specifically for negation and fine-grained composition.
CoN-CLIP and NegationCLIP improve sensitivity to negated concepts, while CoLLM uses LLM-generated triplets for composed image retrieval~\cite{Singh_2025_WACV,Park_2025_ICCV,Huynh_2025_CVPR}.

% Negation is particularly challenging for embedding-based retrieval.
% Collectively, these findings expose a limitation of representations dominated by aggregate semantic similarity: candidates may appear highly similar to a query despite violating one or more of its essential constraints.
% Our objective directly supervises this distinction by organizing candidates according to their degree of constraint satisfaction.

\paragraph{Synthetic supervision for retrieval.}
Large language models have increasingly been used to generate targeted supervision for dense retrieval.
Syntriever and SyNeg generate queries, relevant examples, or challenging negatives that expose distinctions absent from conventional training data~\cite{kim-baek-2025-syntriever,li2024syneg}.
Promptagator \citep{dai2022promptagatorfewshotdenseretrieval} uses LLMs for few-shot dense retrieval.
\citet{wang2024improvingtextembeddingslarge, lee2024geckoversatiletextembeddings} distill embeddings from LLMs.
% \red{E5-mistral synthetic data (https://arxiv.org/abs/2401.00368), Gecko (https://arxiv.org/abs/2403.20327)}
In product retrieval, Agapaki and Gill combine structured product information with LLM-based relevance assessment to identify difficult training examples~\cite{agapaki2026negative}.
% These approaches show that synthetic data can control which distinctions a retriever encounters during training.
% Our use of generated information differs in an important respect: generated feature annotations specify the candidate structure, but do not directly determine the supervision target.
% Given these annotations and the query constraints, the target distance is computed deterministically as a discrete feature-edit distance.

% \paragraph{Multimodal compositional retrieval.}
% Compositional constraints create analogous difficulties in vision--language retrieval.
% CLIP learns a shared embedding space for images and text, whereas the CIRR benchmark and methods such as CLIP4CIR study retrieval from queries that combine a reference image with a textual modification~\cite{pmlr-v139-radford21a,Liu_2021_ICCV,Baldrati_2022_CVPR}.
% More recent approaches introduce supervision specifically for negation and fine-grained composition.
% CoN-CLIP and NegationCLIP improve sensitivity to negated concepts, while CoLLM uses LLM-generated triplets for composed image retrieval~\cite{Singh_2025_WACV,Park_2025_ICCV,Huynh_2025_CVPR}.
% This literature further demonstrates that general semantic alignment does not reliably encode whether each positive and negative constraint is satisfied.

% \paragraph{Positioning of our work.}
% Prior research has improved compositional retrieval through hard-negative mining, teacher-derived graded supervision, Boolean-aware modeling, attribute-sensitive representations, and synthetic example generation.
% Our work connects these directions by introducing a discrete, query-relative structural distance: the minimum number of query-relevant feature edits required for a candidate to satisfy the query.
% Unlike teacher scores or learned relevance estimates, this distance is deterministic once the query constraints and candidate features are specified.
% We use it as graded supervision for first-stage retrieval, encouraging distances in the learned embedding space to reflect degree of constraint satisfaction.
% We evaluate the resulting formulation in both textual and image-based product retrieval.






\section{Dataset Construction}
We seek to evaluate models on a test dataset containing query-document pairs where the document satisfies all constraints in the query.
% At the same time, we want the training dataset to contain hard negative examples for those queries that do not satisfy all their constraints.
To obtain fine-grained attributes tailored to each document pair, we construct a query synthesis pipeline based on the observation that current generative LLMs can extract the necessary attributes.
In effect, this pipeline seeks to distill fine-grained attribute extraction from large, generative models into small embedding models.
\begin{figure}[t]
    \centering
    \includegraphics[width=\linewidth]{figs/dataset-creation.pdf}
    \caption{Beginning with documents 1) from the ESCI Amazon products on In-Shop apparel images datasets, we choose a positive product and a similar hard negative product based on dataset metadata. We choose an easy negative at random. 2) We extract attributes that apply to one, both, or neither product. 3) We randomly select attributes, using them to construct a query that fully applies to the positive example by negating attributes that do not apply to it. We also calculate the number of said attributes $d$ that do not apply to the hard negative. We use an LLM to rephrase the query to reduce key phrase overlap between the product and the query, producing 4) the final dataset.}
    \label{fig:dataset-creation}
\end{figure}



We target two domains, e-commerce products and images, to demonstrate the multimodal capabilities of our system.
For e-commerce, we choose the ESCI corpus of Amazon products \citep{reddy2022shopping} due to its large size.
For images, we choose the DeepFashion In-Shop dataset because it contains images with subtle but objectively differing attributes.

\begin{table}[t]
\centering
\vspace{1em}
\small
\input{figs/dataset_summary_table.tex}
\caption{Dataset summary statistics. Document counts for the Main dataset include hard negative pairs while Human documents are positive-only.}
\label{tab:datasets}
\end{table}


\paragraph{Structured attribute selection}
To extract document attributes, we prompt Gemini 2.5 Flash with a matched positive and negative document pair, using the prompt of Appendix~\ref{prompt:text-features} for products and Appendix~\ref{prompt:image-features} for images.
\darkgreen{We select candidates for hard negatives in ESCI using the pre-existing Substitute and Irrelevant labels, retaining only original-query/positive-product groups with at least two distinct eligible candidates across these labels.}
\darkgreen{We apply this requirement before attribute generation and recheck it after generation and split filtering.}
\darkgreen{For DeepFashion, we select a random example from the same category, such as ``blouses'' or ``graphic tees'', excluding all views and colorways of the positive garment.}
In the DeepFashion dataset, we include images and descriptions of both products in the prompt, instructing the model not to generate attributes that are not visible in the images.

We instruct it to generate up to 5 attributes that apply to only the first example, only the second example, to both examples, or to neither example, a pool of up to 20 attributes.
We then select a random subset of these attributes to compose the query, negating attributes that do not apply to the positive example, ensuring that at least one attribute is not satisfied by the hard negative example.
Oversampling here is useful to ensure that selected attributes are not merely the most salient.
By conditioning attribute generation on the negative document, similar to \citet{mayfield2023syntheticcrosslanguageinformationretrieval, zeng2026dualviewtraininginstructionfollowinginformation, vedantam2017contextawarecaptionscontextagnosticsupervision} we ensure that forbidden attributes could plausibly apply to a document similar to the positive.
% Thus, we obtain a set of $k$ attributes, all of which are satisfied by the positive example and $k'$ of which are \emph{not} satisfied by the hard negative example, where $1 \leq k' \leq \min(k, 10)$. % important! add back in below somewhere

We construct a raw query with the template ``I am looking for: \{item\} that has: \{positive\} + \{common\} attributes; and does not have: \{negative\} + \{neither\} attributes''.
We then prompt Gemini 2.5 Flash to rewrite the raw query based on five human-written in-context examples, replacing attribute phrases with synonyms (Appendix~\ref{prompt:rephrase}); the prompt includes the positive document so synonyms stay consistent with it, requires at least half of the attributes to be reworded, and forbids changing numeric values.
Rewriting the query in this way is designed to improve the realism and diversity of queries and reduce surface form overlap between queries and positive documents.
This yields a (query, positive) pairs.

\paragraph{Baseline negative sampling}
\green{todo: refresh dataset statistics and experimental results for the candidates2 text dataset and metadata-matched Baseline.}
\darkgreen{Our Baseline samples comparison negatives using the same source metadata as \textsc{gold} mining.}
\darkgreen{For text, we sample uniformly from distinct eligible Substitute or Irrelevant products associated with the same original ESCI query and positive product in the training split.}
\darkgreen{For images, we sample uniformly from training images in the positive's category, excluding all views and colorways of its garment.}
\darkgreen{The comparison document used during query generation remains eligible for sampling.}
\darkgreen{All strategies retain the same additional easy-negative triples: split-local random products for text and images from a different category for DeepFashion.}
\darkgreen{Only the training comparison negatives change across strategies; validation and test data are shared.}
% We also compare to the state-of-the-art hard negative mining algorithm NV-Retriever \citep{moreira2024nv}, described in \red{...}
\darkgreen{We also compare to an adaptation of NV-Retriever's positive-aware hard negative mining strategy \citep{moreira2024nv}, a key component of NV-Embed-v2 \citep{lee2024nv}, the highest-scoring retrieval model in Table~5 of \citet{wischounig2026negative}.}
Regardless of the strategy, this produces (query, positive, negative) triples.

\paragraph{\textsc{gold} mining}
We note that conditioning query generation on random documents as described above makes these documents, by construction, hard negatives with respect to the query.
\green{todo: check for correctness after baseline dataset update.}
Specifically, each document will satisfy exactly distance $d$ conditions in the query, where $1 \leq d \leq \min(k, 10)$.
We term this strategy \textsc{gold} mining.
Thus, \textsc{gold} mining produces (query, positive, negative, distance) tuples, where $d$ is consumed by the \textsc{gold} mining objective described in Section \red{tbd}.
Although \textsc{gold} mining here occurs in parallel with query generation, in principle this strategy is dataset-agnostic.
The strategy could be performed in reverse by pairing existing conjunctive queries with negative documents and calculating $d$. 


% We also sample an easy negative at random from the same split.
% The easy negative is drawn from a different category than the positive for DeepFashion.
% This yields a matching (query, positive, easy negative) triple. 
% We divide this dataset into train, validation, and test splits, ensuring that no query or document appears in multiple splits.
% Note that the hard negative mining procedure can be applied to virtually any dataset or modality on which an LLM can perform binary QA.
% Thus, we consider it the first contribution of this paper.

\paragraph{Human labeled queries}

Annotators manually annotated 133 and 83 text and image queries, respectively, based on positive documents.
We use these query-document pairs to validate that models trained on Main data generalize to human-generated queries.


\section{Methods}
\label{sec:methods}

Numerous training objectives have been proposed for dense retrieval.
We choose Margin-MSE \citep{hofstatter2020improving}, CoSENT \citep{huang2024cosent}, SigLIP \citep{tschannen2025siglip}, and InfoNCE \citep{oord2019representationlearningcontrastivepredictive} as our baseline loss functions based on their applicability to triple-supervised training and coverage of high-level loss strategy families.
Margin-MSE regresses query-document pairs onto a fixed score different; CoSENT ranks pairs ordinally; SigLIP classifies in-batch pairs independently; InfoNCE scores each in-batch candidate against the others using a softmax function.
We include their definitions in Appendix \ref{app:loss-functions} for reference.
Because batch size affects CoSENT, SigLIP, and InfoNCE training dynamics through pairwise in-batch comparisons, we hold batch size constant across all runs.

\subsection{Baseline Details}

For InfoNCE, SigLIP, and batched Margin-MSE, each of the $B$ queries is scored against all $2B$ documents in the batch, yielding $2B^2$ query--document scores.
Documents from other triples are treated as easy negatives; each query's own paired negative may be hard or easy, chosen uniformly at random.
InfoNCE uses one softmax over all candidates.
Margin-MSE and SigLIP weight the mean own-row and cross-row losses equally: the own-row group contains the paired negative for Margin-MSE, and both the positive and paired negative for SigLIP.
CoSENT compares scores of supplied query--document pairs, while pairwise MSE regresses each pair's score independently.
Within each loss family, candidate construction is shared across negative-selection strategies.
We design our experiments such that each model sees each training triple once per epoch.
We train every model for one epoch with 32 rows per forward pass and 4 forward passes per optimizer step.
Pairwise MSE emits two pairs per triple.
% We train baseline models using hyperparameters listed in Section \ref{sec:model-details}, concluding that InfoNCE is the strongest.
To study the NV-Retriever baseline, we perform a seach across NV-Retriever mining paramters independently for each loss function, described in Section \red{??}.

% \red{To further strengthen our baseline, we adopt the positive-aware hard negative mining of NV-Retriever \citep{moreira2024nv}, which the survey of \citet{wischounig2026negative} classifies under false-negative filtering and which underlies the top retrieval model on the MTEB leaderboard in their Table 5 \textcolor{orange}{(cite NV-Embed-v2)}, and apply it to InfoNCE.}
% We choose the best of $n=8$ hyperparameter combinations based on a grid search \red{(Appendix \ref{??})}.


% \begin{table}[t]
% \centering
% \input{figs/baseline_loss_table.tex}
% \caption{Validation split recall of baseline loss functions trained on binary labels with randomly selected negatives, $n=3$. +NV indicates the model was trained using negatives mined with NV-Retriever.}
% \label{tab:baseline-loss}
% \end{table}


\subsection{\textsc{gold} Reweighting}


Traditional Margin-MSE loss \citep{hofstatter2020improving} is defined as

% \begin{equation}
%     \mathcal{L}{q, x^+, x^-} = \Big( \inner{f(q)}{f(x^+)} - \inner{f(q)}{f(x^-)} \Big)^2
% \end{equation}
\begin{equation}
    \label{eq:trad-mse}
    \mathcal{L}(q, x^+, x^-) = \Big( \inner{f(q)}{f(x^+)} - \inner{f(q)}{f(x^-)} - y \Big)^2 .
\end{equation}

Here, $q$ is a query, $x^+$ is a relevant document, $x^-$ is a less relevant document, and $y$ is a scalar representing the margin between their relevance as determined by a teacher model.
$f$ is a function mapping queries or documents to their embeddings. 
For loss functions we choose cosine similarity as our inner product.

% Retrieval scores use cosine similarity, $\operatorname{cos}(a,b)=a^\top b/(\|a\|\|b\|)$, while $\inner{a}{b}=a^\top b$ denotes the Euclidean inner product.
% Because $f$ returns unit-norm embeddings, $\operatorname{cos}(f(q),f(x))=\inner{f(q)}{f(x)}$, so a difference of cosine scores equals $\inner{f(q)}{f(x^+)-f(x^-)}$.
% The difference vector and the probe vectors below are not normalized.



Formally, following \citet{lu2025multiconir}, let query $q_k$ be a query composed of $k$ conditions $C=\{c_1, c_2, ..., c_k\}$ with $k \in \{1, ..., K\}$.
We say a document $x$ is a positive document, $x^+$, with respect to $q_k$ if $x^+ \models C$, i.e., $x^+$ satisfies all $k$ conditions.
Otherwise, we say the document $x^-$ is negative with respect to $q_k$.
$x^-_h$ and $x^-_e$ denote hard and easy negatives, respectively.




% We introduce Attribute Hamming Loss, which is a form of Margin-MSE loss designed to distil knowledge from generative models rather than teacher embedding models.
We introduce \textsc{gold} reweighting, a strategy for assigning partial credit to hard negative examples.
The idea behind \textsc{gold} reweighting is to modify a loss function such that loss of a query and a document is proportional to the number of constraints violated by a document.
Because CoSENT relies on ordering for loss, \textsc{gold} reweighting on CoSENT instead treats all hard negatives as rank 1/2.
Formally, we define attribute Hamming distance $d$ as
% Attribute Hamming Loss replaces the target $y$ with an integer representing the number of constraints in the query violated by $x^-$, rescaled by a constant $V$.
% Training operates on triples $(q_k, x^+, x^-)$ in which the positive satisfies every condition of the query, and the predicted margin between the positive and negative similarities is regressed onto the number of query constraints the negative violates:

\begin{equation}
    d(x, q) =
    % \begin{cases}
        % \eta & \text{if } x \text{ is an easy negative} \\
        \sum_{i=1}^{k} \ind{x \not\models c_i} %& \text{otherwise}
    % \end{cases}
\end{equation}

We then modify each loss function, with the modification marked by an underbrace:

\paragraph{Margin-MSE:}
We set the margin term to $d/V$ and apply the loss across all query--document pairs in the batch:

\begin{equation}
\label{eq:hamming}
    \mathcal{L_\text{MSE}}(q_k, x^+, x^-) = \Big( \inner{f(q_k)}{f(x^+)} - \inner{f(q_k)}{f(x^-)} - \underbrace{\tfrac{d}{V}}_{\text{ours}} \Big)^2
\end{equation}

The violation scale hyperparameter $V$ controls the penalty per violation.
Let $K^*$ be the maximum number of constraints violated by any example.
In Margin-MSE only, we set $d(x,q)=\eta \geq K^*$ when $x$ is an easy negative.
Choosing $V \geq \text{max}(\eta, K^*)$ ensures the $d/V$ term is in $[0, 1]$.

\paragraph{InfoNCE:}
Every document in a batch of $B$ triples is a candidate for every query: the $B$ positives and the $B$ hard negatives, the set $X$ of $2B$ candidates.
InfoNCE trains the softmax over $X$ toward a one-hot target on the query's own positive. We keep the logits and change the target: the own positive keeps mass $1$, the own hard negative receives mass $e^{-s\,d/V}$ for its measured $d$, and every candidate drawn from another row, positive or negative, is a random product with respect to $q_i$ and receives no mass. Rows are normalized to sum to one:

\begin{align}
    \label{eq:ours-infonce-target}
    T_i(x) &=
    \begin{cases}
        1 & x = x_i^+ \\
        e^{-s\, d(x_i^-, q_i) / V} & x = x_{h,i}^-\\\
        0 & \text{otherwise,}
    \end{cases} \\[0.5em]
    \label{eq:ours-infonce}
    \mathcal{L}_\text{INCE}(q_i) &= - \sum_{x \in X} \underbrace{\frac{T_i(x)}{\sum_{x' \in X} T_i(x')}}_{\text{ours}} \log \frac{\exp\big( s \inner{f(q_i)}{f(x)} \big)}{\sum_{x' \in X} \exp\big( s \inner{f(q_i)}{f(x')} \big)} .
\end{align}

As $V \rightarrow 0$, this becomes equivalent to the original InfoNCE loss.

\paragraph{SigLIP:}
In place of the original binary target, we weight each pair with $T_i(x)$:

\begin{equation}
    \label{eq:ours-siglip}
    \mathcal{L}_\text{SigLIP} = - \frac{1}{B} \sum_{i=1}^{B} \sum_{x \in X} \left[ \underbrace{T_i(x)}_{\text{ours}} \log \frac{1}{1 + e^{-\left( s \inner{f(q_i)}{f(x)} + b \right)}} + \underbrace{\big( 1 - T_i(x) \big)}_{\text{ours}} \log \frac{1}{1 + e^{s \inner{f(q_i)}{f(x)} + b}} \right] ,
\end{equation}


\paragraph{CoSENT:}
We assign rank 1 to positives, 1/2 to hard negatives, and 0 to easy negatives. 
The algorithm otherwise remains unchanged.
% \red{CoSENT reads its labels only through their ordering, so we assign each pair $(q_i, x_i)$ one of three ranks:}

% The unifying factor among these modifications is to regress towards the condition 


The unifying factor among these modifications is to assign partial credit to hard negatives, compared to easy negatives.
For InfoNCE and Margin-MSE, this regresses towards the condition

\begin{equation}
    \inner{f(q)}{f(x_1)} - \inner{f(q)}{f(x_2)} = (d(q, x_2) - d(q, x_1))/V
\end{equation}





\subsection{{Relation to Affine Probing}}

\begin{figure}
    \centering
    \includegraphics[width=0.8\linewidth]{figs/embedding_fig.pdf}
    \caption{An idealized embedding space in $\mathbb{R}^4$, drawn by dropping the fourth coordinate.
    Three documents, each with exactly one attribute, have mutually orthogonal visible parts and put the rest of their unit length into a fourth component the query cannot see.
    With $V = 3/\sqrt{2}$ as the violation scale, the probe sum has norm $\|b_{\mathrm{red}} + b_{\mathrm{zipper}}\| = V$, and $V \inner{f(q)}{f(x)}$ counts the conditions $x$ satisfies: $x_{\mathrm{red}}$ and $x_{\mathrm{zipper}}$ score $1/V = \sqrt{2}/3$, $x_{\mathrm{cotton}}$ scores $0$, and the positive $x_{\mathrm{red,zipper}}$, which the probes force to be the visible sum of $x_{\mathrm{red}}$ and $x_{\mathrm{zipper}}$, scores $2/V$.
    Any larger $V$ is absorbed by the document embeddings moving more of their length into the fourth component the query cannot see.}
    \label{fig:embedding}
\end{figure}




% For a given $c_i$, 
Following \citet{garg2026many}, we will assume that inputs come from some set $L = Q \cup D$ consisting of queries or documents composed of text or images.
Let $f : L \rightarrow \mathbb{R}^d$ give the embeddings of dimension $d$.
We let $z_i : L \rightarrow \{0,1\}$ be an attribute of a query or corresponding product, for example, whether a query requires a red product, or whether the product in question is red. 
Note that the value of $z_i$ depends on the underlying input, not its embedding.


% Under the Linear Representation Hypothesis, we assume that for a set of features $z_1, ..., z_m : L \rightarrow \mathbb{R}$

For a document set $S \subseteq D$ and tolerance $\epsilon>0$, we say a feature is $(\epsilon, S)$-affinely recovered by a probe $(b_i,\beta_i)\in\mathbb{R}^d\times\mathbb{R}$ if
% Let $b_i \in \mathbb{R}^d$ be a probe vector for attribute $z_i$ such that

\begin{equation}
    |\inner{b_i}{f(\ell)} + \beta_i - z_i(\ell)| < \epsilon
\end{equation}

for all $\ell \in S$.

Under the assumptions of Theorem~\ref{thm:forward}, affine recoverability of attributes permits low \textsc{gold}-reweighted Margin-MSE loss.
Conversely, uniformly low loss across all positive--negative comparisons for two queries differing by one condition, each with nonempty positive and negative sets, implies approximate affine recovery of that condition's attribute (Theorem~\ref{thm:back}).
Binary-target Margin-MSE provides no such guarantee: zero loss can coexist with an unrecoverable attribute (Proposition~\ref{prop:binary}).
Finally, for documents with at most $k'$ active attributes from a vocabulary of size $m$, queries with at most $k$ conditions, and $V\ge kk'$, we construct unit document and query embeddings in dimension $d=O\!\left(\frac{k'^2}{\epsilon^2}\log m\right)$ that affinely recover every attribute with error below $\epsilon$ and achieve loss below $(2k\epsilon/V)^2$ on every valid comparison (Corollary~\ref{cor:capacity}).


\subsection{Implementation and Evaluation Details}

We fine-tune \texttt{all-mpnet-base-v2} on text and \texttt{clip-ViT-B-32} on images, both as distributed by Sentence-Transformers \citep{reimers-gurevych-2019-sentence}, normalizing every representation to unit length.
We optimize with AdamW \citep{loshchilov2017decoupled} using a learning rate $2 \times 10^{-5}$, a cosine schedule warming up over the first 10\% of steps, and bf16, training each condition three times and reporting the mean over those trials.
Training details are identical for baseline and \textsc{gold} models.

For \textsc{gold} models, we sweep $V \in \{10,20,40,80\}$ for InfoNCE and SigLIP, and $V \in \{20,40,80\}$ with $\eta \in \{10,20,40\}$ for Margin-MSE, selecting settings on the validation set.
The CoSENT modification introduces no additional hyperparameters.
We perform ablations that only use our dataset creation process without modifying baseline loss functions.

% \subsection{Evaluation Details}

We choose recall@20 and recall@5 as our primary metrics for text and image retrieval, respectively.
For text retrieval, we report Recall@5 because this more stringent cutoff provides greater discrimination among high-performing strategies.

\section{Results and Discussion}

\begin{figure}[t]
    \centering
    \begin{subfigure}{\linewidth}
        \centering
        \includegraphics[width=\linewidth,trim=0 11bp 0 0,clip]{figs/text_vs_image_recall.pdf}
        \phantomcaption
        \label{fig:text-vs-image-recall}
    \end{subfigure}\par
    \vspace{-1em}
    \nointerlineskip
    \begin{subfigure}{\linewidth}
        \centering
        \includegraphics[width=\linewidth,trim=0 0 0 3bp,clip]{figs/recall_vs_winrate.pdf}
        \phantomcaption
        \label{fig:recall-vs-winrate}
    \end{subfigure}
    \caption{(a) Retrieval recall and hard-negative win rate on the test split.
    Arrows indicate that one training strategy builds on another.
    Unfilled shapes indicate baselines.
    Recall versus positive/hard-negative win rate for text (b) and images (c), $n=3$ per point.}
    \label{fig:text-vs-image}
\end{figure}


We present results in Figure \ref{fig:text-vs-image}.
On the larger Main test set, we see significant improvement compared to baselines when applying both \textsc{gold} mining and \textsc{gold} reweighting.
Overall accuracy is stratified by loss function, with InfoNCE ranking highest on both modalities, both for our proposed methods and baselines.
Both the NV-Retriever baseline and \textsc{gold} mining significantly increase recall on text, but improvement on images is mixed and modality dependent.
\textsc{gold} mining is more effective than NV-Retriever on both modalities using the two strongest loss functions, InfoNCE and SigLIP.
After applying \textsc{gold} reweighting, our method surpasses both random and NV-mined baselines in seven of eight loss function--modality pairs on Main dataset queries.
For InfoNCE on images, \textsc{gold} mining slightly outperforms both baselines, whereas mining plus reweighting falls short of the strongest baseline by approximately 0.001 recall.
% \darkgreen{The exception is InfoNCE on images, where our method falls short of the strongest baseline, random negatives, by approximately 0.001 recall.}

We observe that across 7/8 loss function--modality pairs, \textsc{gold} mining increases both recall and win rate compared to the random baseline.
Applying \textsc{gold} reweighting improves synthetic recall in six of eight loss function--modality pairs while consistently reducing hard-negative win rate.
That is, \textsc{gold} mining improves accuracy more on hard negatives while \textsc{gold} reweighting trades accuracy on hard negatives for accuracy on easy negatives.
This is the expected result of penalizing errors in hard negatives less than those in easy negatives in the loss function.
Thus, the application of \textsc{gold} reweighting, and in particular the $V$ scale, should be considered in the context of the dataset.
Datasets containing more hard negatives per query may benefit more from \textsc{gold} mining.
Applications requiring maximum recall will benefit more from 

We see in Figure~\ref{fig:infonce-condition-counts} that InfoNCE accuracy for both \textsc{gold} mined + reweighted and the random baseline is positively correlated with the number of positive (required) attributes in both modalities.
However, accuracy appears to be largely uncorrelated with the number of negative (forbidden) attributes.
The correlation with required attributes in conjunctive queries is unsurprising given that keyword or synonym overlap between queries and positive documents is tightly correlated with the number of conditional queries.
The lack of correlation with forbidden attributes suggests that overall, embedding models neither induce false positives out of documents containing those attributes nor contribute to magnifying positive documents in the way that required attributes do.
Other loss functions show similar behavior.

To validate the use of synthetic queries in the Main dataset, we compare model accuracy on human-generated versus synthetic queries.
Across all 96 models, we find a strong linear relationship between recall on human versus synthetic queries (Pearson $r \geq 0.83$, $R^2 \geq 0.70$).
At the same time, the correlation in model ordering ranges from moderately positive (Spearman $\rho = 0.47$) for text to strongly positive (\darkgreen{$\rho = 0.80$}) for images.
Surprisingly, despite being trained on LLM-generated queries, average model recall on human-generated queries exceeds that on synthetic queries.
This cannot be explained simply by the lower average number of conditions in human queries because condition count is in fact positively correlated with recall (Figure~\ref{fig:infonce-condition-counts}).
Human queries also reuse more wording from their positive text documents than paired synthetic queries (unigram precision 58.6\% vs.\ 40.4\%; bigram precision 22.4\% vs.\ 10.6\%, respectively).
Furthermore, human queries are shorter on average than synthetic queries (Table~\ref{tab:datasets}).



\begin{figure}[t]
    \centering
    \begin{minipage}[c]{0.355\linewidth}
        \includegraphics[width=\linewidth]{figs/human_vs_synthetic_text_recall.pdf}
    \end{minipage}\hfill
    \begin{minipage}[c]{0.355\linewidth}
        \includegraphics[width=\linewidth]{figs/human_vs_synthetic_image_recall.pdf}
    \end{minipage}\hfill
    \begin{minipage}[c]{0.27\linewidth}
        \centering
        \input{figs/human_synthetic_correlations.tex}
    \end{minipage}
    \caption{Main (synthetic) versus human recall for text (Recall@5, left) and images (Recall@20, center), with correlation statistics.
    Each plot includes all 48 model sseds; colors and markers follow Figure~\ref{fig:text-vs-image-recall}, and dashed lines show linear fits.}
    \label{fig:human-synthetic-correlation}
\end{figure}

\begin{figure}[t]
    \centering
    \includegraphics[width=\linewidth]{figs/infonce_condition_counts.pdf}
    \caption{InfoNCE by positive and negative query condition count showing correlation with positive conditions but not negative conditions. Other loss functions show identical trends.}
    \label{fig:infonce-condition-counts}
\end{figure}

\section{Conclusion}

We present \textsc{gold} mining, a multimodal hard negative mining algorithm that improves recall on conjunctive queries.
We also present \textsc{gold} reweighting, a loss function modification strategy that, when combined with \textsc{gold} mining, outperforms random and NV-mined baselines in seven of eight Main evaluation settings across four loss functions and two modalities, with a deficit of approximately 0.001 recall in the remaining setting.
Theoretical analysis of \textsc{gold} reweighting applied to Margin-MSE shows that under suitable query and comparison coverage assumptions, low loss implies approximate affine recoverability of attributes, a property not guaranteed to be satisfiable by the standard objective. \red{add a note about generalizability}.

% Keep the results figures within Results and Discussion.
\FloatBarrier




\bibliography{related-work}
\bibliographystyle{iclr2026_conference}

\appendix

\section{Theoretical Analysis}


We analyze the hard-negative Margin-MSE objective with measured violation counts, assuming the presence of no easy negatives.
Theorem~\ref{thm:forward} constructs cosine queries from affine attribute probes under explicit norm and scale assumptions.
Theorem~\ref{thm:back} shows that, under complete comparison coverage, low count loss on a query pair differing in one condition implies affine recovery of the deleted attribute.
Proposition~\ref{prop:binary} compares this result with binary supervision within Margin-MSE, which can fit the same query pair without preserving the deleted attribute.
These results describe feasible representations and information required by low-loss embeddings.
Corollary~\ref{cor:capacity} gives a unit-norm construction with logarithmic dependence on the attribute vocabulary size under an explicit sufficient bound on $V$.

\subsection{Setup}

Let $S$ be a document set, let $f$ map queries and documents to unit vectors in $\mathbb{R}^d$, and fix $V>0$.
For a condition $c$, define its satisfaction indicator $h_c(x)=\ind{x\models c}$, equal to $z_i(x)$ when $c$ requires attribute $i$ and $1-z_i(x)$ when it forbids that attribute.
For a query $q_C$ with $k=|C|\ge1$ conditions, define
\begin{equation}
\label{eq:app-label}
    H_C(x)=\sum_{c\in C}h_c(x),
    \qquad y_C(x)=k-H_C(x).
\end{equation}
Write $P_C=\{x\in S:H_C(x)=k\}$ and $N_C=S\setminus P_C$.
For $p\in P_C$ and $n\in N_C$, the target count is $y_C(n)=H_C(p)-H_C(n)$ and the cosine margin is
\begin{equation}
\label{eq:app-obj}
\begin{aligned}
    M_C(p,n)&=\operatorname{cos}(f(q_C),f(p))-\operatorname{cos}(f(q_C),f(n)) \\
             &=\inner{f(q_C)}{f(p)-f(n)},
    \qquad \mathcal{L}(q_C,p,n)=\big(M_C(p,n)-y_C(n)/V\big)^2.
\end{aligned}
\end{equation}
The inner product is Euclidean; the equality uses unit embeddings and does not normalize the difference $f(p)-f(n)$.

An affine probe $(b,\beta)$ recovers a function $h$ with error $\epsilon>0$ on $S$ if
\begin{equation}
\label{eq:app-rec}
    |\inner{b}{f(x)}+\beta-h(x)|<\epsilon
    \qquad\text{for all }x\in S.
\end{equation}
If $(b_i,\beta_i)$ recovers $z_i$, then $(-b_i,1-\beta_i)$ recovers $1-z_i$ with the same error, so recovering a condition indicator and recovering its attribute are equivalent.

\subsection{Main Results}

\begin{theorem}[Affine recoverability permits low loss]
\label{thm:forward}
Let $U=\operatorname{span}\{f(x):x\in S\}$ be a proper subspace of $\mathbb{R}^d$.
Suppose every $h_c$, $c\in C$, has an affine probe $(b_c,\beta_c)$ with error $\epsilon>0$ on $S$ and $b_c\in U$, and put $u=\sum_{c\in C}b_c$.
Using the attribute probes above, $u$ is their signed sum: required attributes contribute $b_i$ and forbidden attributes contribute $-b_i$.
If $V\ge\|u\|$, the unit query
\begin{equation}
\label{eq:app-unitq}
    f(q_C)=(u+w)/V,
    \qquad w\perp U,
    \qquad \|w\|^2=V^2-\|u\|^2
\end{equation}
attains $\mathcal{L}(q_C,p,n)<(2k\epsilon/V)^2$ for every $(p,n)\in P_C\times N_C$.
\end{theorem}

\begin{proof}
The assumptions ensure that $w$ exists, that $\|f(q_C)\|=1$, and that $w$ contributes nothing to document scores.
Set $r_c(x)=\inner{b_c}{f(x)}+\beta_c-h_c(x)$.
The probe biases cancel in each difference, giving
\begin{equation}
    VM_C(p,n)-y_C(n)=\sum_{c\in C}\big(r_c(p)-r_c(n)\big).
\end{equation}
Each summand has magnitude below $2\epsilon$, so $|M_C-y_C/V|<2k\epsilon/V$; squaring proves the claim.
\end{proof}

Corollary~\ref{cor:capacity} below constructs document embeddings and probes satisfying these assumptions when $V\ge kk'$, where $k'$ bounds the number of active attributes per document.

\begin{theorem}[Low loss implies affine recoverability]
\label{thm:back}
Let $\delta>0$, suppose $P_C,N_C$ are nonempty, and assume $\mathcal{L}(q_C,p,n)<\delta^2$ for every $(p,n)\in P_C\times N_C$.
Choose $x_0\in P_C$ and define
\begin{equation}
\label{eq:app-agg}
    a_C=k-V\inner{f(q_C)}{f(x_0)},
    \qquad G_C(x)=V\inner{f(q_C)}{f(x)}+a_C.
\end{equation}
Then $|G_C(x)-H_C(x)|<2V\delta$ for every $x\in S$.
For $c\in C$, if $D=C\setminus\{c\}$ is also queried and satisfies the same hypotheses on $S$ with the same $V,\delta$, define $G_D$ using $|D|$ in place of $k$ and the same reference $x_0\in P_C\subseteq P_D$.
Then $G_C-G_D$ affinely recovers $h_c$, and hence its underlying attribute, with error $4V\delta$; the probe for $h_c$ is
\begin{equation}
\label{eq:app-probe}
    b=V\big(f(q_C)-f(q_D)\big),
    \qquad \beta=a_C-a_D.
\end{equation}
\end{theorem}

\begin{proof}
For $x\in N_C$, $G_C(x)-H_C(x)=y_C(x)-VM_C(x_0,x)$ has magnitude below $V\delta$ by the loss assumption.
For $x\in P_C$, choose $n\in N_C$ and subtract the controlled margins $M_C(x,n)$ and $M_C(x_0,n)$.
Their targets agree, giving $|G_C(x)-H_C(x)|=V|M_C(x,n)-M_C(x_0,n)|<2V\delta$.
Applying the same argument to $D$ and using $H_C-H_D=h_c$ gives
\begin{equation}
    |(G_C-G_D)(x)-h_c(x)|
    \le |G_C(x)-H_C(x)|+|G_D(x)-H_D(x)|<4V\delta.
\end{equation}
\end{proof}

At fixed $\delta$, the recovery bound $4V\delta$ grows linearly with $V$; zero loss gives exact recovery.

To compare binary and count supervision within Margin-MSE, we use the same queries and comparisons with binary target $1$ for every positive--negative pair.

\begin{proposition}[Binary Margin-MSE can discard attributes]
\label{prop:binary}
There are unit embeddings and a query pair $C=\{c_1,c_2\}$, $D=\{c_1\}$, both conditions positive, with zero binary Margin-MSE loss
\begin{equation}
    \mathcal{L}_{\mathrm{bin}}(q_C,p,n)
    =\big(M_C(p,n)-1\big)^2
\end{equation}
on every positive--negative pair for both queries, although every affine probe for $z_2$ has maximum error at least $1/2$ on the document set.
\end{proposition}

\begin{proof}
Take one document for each $(z_1,z_2)\in\{0,1\}^2$ and choose $f(q_C)=e_1$, $f(q_D)=e_2$ in $\mathbb{R}^3$, with
\begin{equation}
    \begin{array}{c|c}
        (z_1,z_2)&f(x)\\ \hline
        (0,0)&(-\tfrac12,-\tfrac12,\tfrac{1}{\sqrt2})\\[2pt]
        (0,1)&(-\tfrac12,-\tfrac12,\tfrac{1}{\sqrt2})\\[2pt]
        (1,0)&(-\tfrac12,\tfrac12,\tfrac{1}{\sqrt2})\\[2pt]
        (1,1)&(\tfrac12,\tfrac12,\tfrac{1}{\sqrt2}).
    \end{array}
\end{equation}
All vectors have unit norm and give margin $1$ on every required pair.
Any probe assigns the same value to $(0,0)$ and $(0,1)$, so one of its errors for $z_2$ is at least $1/2$.
\end{proof}

Using our \textsc{gold} reweighting, on the other hand, $(0,0)$ and $(0,1)$ have different targets because they violate two conditions and one condition of $C$, respectively.
Per Theorem~\ref{thm:back}, sufficiently small loss on these two queries implies that $z_2$ is approximately recoverable by an affine probe, even without a query for $z_2$ alone.
% Theorem~\ref{thm:back} shows that uniformly small count loss on these two queries makes $z_2$ approximately recoverable, even without a query for $z_2$ alone.

\subsection{Capacity with Unit Embeddings}

We adapt the approximately orthogonal construction of \citet{garg2026many} by applying one common scale and adding two coordinates.
One coordinate normalizes the document embeddings; the other normalizes the query embeddings.

\begin{corollary}[Capacity for cosine Margin-MSE]
\label{cor:capacity}
Let $m\ge2$, $k\ge1$, and $1\le k'\le m$ be integers, let $0<\epsilon\le1$, and suppose $z(x)\in\{0,1\}^m$ has at most $k'$ nonzero entries for every $x\in S$.
For every $V\ge kk'$, there exist unit document and query embeddings in dimension
\begin{equation}
\label{eq:app-capacity}
    d=O\!\left(\frac{k'^2}{\epsilon^2}\log m\right)
\end{equation}
such that every attribute is affinely recovered with error $\epsilon$ on $S$, and every query with $1\le|C|\le k$, allowing positive and negated conditions, has $\mathcal{L}(q_C,p,n)<(2k\epsilon/V)^2$ for every $(p,n)\in P_C\times N_C$.
\end{corollary}

\begin{proof}
Set $\mu=\epsilon/k'$.
Lemma~5 of the cited work constructs $m$ unit vectors $a_i\in\mathbb{R}^{d_0}$ with $|\inner{a_i}{a_j}|<\mu$ for $i\ne j$, using $d_0=O(\mu^{-2}\log m)$ coordinates.
Here $d_0$ is the dimension before the two normalization coordinates are added; put $A=[a_1\ \cdots\ a_m]$.
Since $\|Az(x)\|\le k'$, define document embeddings and attribute probes in $\mathbb{R}^{d_0+2}$ by
\begin{equation}
\label{eq:app-capacity-padding}
    f(x)=\left(\frac{Az(x)}{k'},\sqrt{1-\frac{\|Az(x)\|^2}{k'^2}},0\right),
    \qquad \widetilde b_i=(k'a_i,0,0).
\end{equation}
Every $f(x)$ has unit norm, including when $z(x)=0$, and the added coordinates contribute nothing to probe outputs.
Since each column has unit norm and distinct columns have inner products bounded in magnitude by $\mu$, we have
\begin{equation}
    \left|\inner{\widetilde b_i}{f(x)}-z_i(x)\right|
    =\left|\sum_{j\ne i}z_j(x)\inner{a_i}{a_j}\right|<k'\mu=\epsilon.
\end{equation}
Project $\widetilde b_i$ onto the document span $U$ to obtain $b_i$; this preserves all probe outputs and gives $\|b_i\|\le k'$.
The last coordinate of every document is zero, so $U$ is a proper subspace.
Positive and negated conditions use $(b_i,0)$ and $(-b_i,1)$, respectively, giving
\begin{equation}
    \left\|\sum_{c\in C}b_c\right\|\le |C|k'\le kk'\le V.
\end{equation}
Theorem~\ref{thm:forward} therefore supplies unit queries with loss below $(2|C|\epsilon/V)^2\le(2k\epsilon/V)^2$.
Substituting $\mu=\epsilon/k'$ into the construction's dimension bound gives $d_0=O(k'^2\epsilon^{-2}\log m)$.
The two added coordinates change this by an additive constant, so $d=d_0+2=O(k'^2\epsilon^{-2}\log m)$.
\end{proof}



\paragraph{Notation.}
\begin{description}\setlength{\itemsep}{0pt}
    \item[$S,f,d$] Document set, unit embedding map, and embedding dimension.
    \item[$z_i,h_c,H_C$] Attribute indicator, condition satisfaction indicator, and number of conditions satisfied.
    \item[$C,k,q_C$] Condition set, its size $k=|C|$, and its query; in Corollary~\ref{cor:capacity}, $k$ bounds the query size.
    \item[$P_C,N_C$] Documents satisfying all conditions and documents violating at least one.
    \item[$b_i,\beta_i$] Affine probe for attribute $z_i$; $(b_c,\beta_c)$ denotes a probe for condition indicator $h_c$.
    \item[$U,u,w$] Document span, sum of condition-probe vectors, and orthogonal query component in Theorem~\ref{thm:forward}.
    \item[$M_C,y_C,V$] Cosine margin, violation count, and fixed positive label scale.
    \item[$\epsilon,\delta$] Uniform probe-error and margin-error tolerances; loss is below $\delta^2$ on each comparison assumed in Theorem~\ref{thm:back}.
    \item[$G_C,a_C$] Affine predictor of $H_C$ and its offset fixed by a reference positive.
    \item[$m,k'$] Attribute vocabulary size and document sparsity in Corollary~\ref{cor:capacity}.
    \item[$A,d_0,\mu$] Matrix of unit attribute columns, dimension before padding, and bound on off-diagonal inner products in the capacity construction.
\end{description}

\section{\red{Full Results}}
\label{app:full-results}

\begin{table}[h]
\centering
\input{figs/main_results_table.tex}
\caption{\red{Test-split recall of every loss and negative-selection group, mean of $n=3$ seeds with the lowest and highest seed in brackets. Bold is the best value in a column; underline is the best within a loss. Synthetic queries are the ones each model trained on; human queries are the human-written set, scored on the same model. Text is Recall@5 and image Recall@20.}}
\label{tab:main-results}
\end{table}

\begin{table}[h]
\centering
\input{figs/significance_table.tex}
\caption{\red{Paired permutation tests for the differences Figure~\ref{fig:text-vs-image} argues about: $A - B$ in recall at each modality's cutoff, with each group represented by the mean of $n=3$ seeds, on the same queries as Table~\ref{tab:main-results}.
The null distribution flips the sign of each positive product's per-query differences (text test queries share products; image and human queries are one per product); $p$ is the one-sided $p$-value for the hypothesis that $A$ is better than $B$, over 10{,}000 Monte Carlo flips with the add-one correction.
Bold is $p < 0.05$.}}
\label{tab:significance}
\end{table}

\section{Loss Functions}
\label{app:loss-functions}

InfoNCE \citep{oord2018representation} treats every document in a batch of $B$ triples as a candidate for every query and is defined as

\begin{equation}
    \label{eq:trad-infonce}
    \mathcal{L}(q_i) = - \log \frac{\exp\big( s \inner{f(q_i)}{f(x_i^+)} \big)}{\sum_{j=1}^{B} \exp\big( s \inner{f(q_i)}{f(x_j^+)} \big) + \exp\big( s \inner{f(q_i)}{f(x_j^-)} \big)} ,
\end{equation}

where $(q_i, x_i^+, x_i^-)$ is the $i$-th triple and the sum runs over the $2B$ positives and negatives of the batch.
The target is one-hot on $x_i^+$, so the loss constrains only the ordering of the candidates.

SigLIP \citep{zhai2023sigmoid} replaces the softmax over candidates with an independent sigmoid classifier for each (query, candidate) pair and is defined as

\begin{equation}
    \label{eq:trad-siglip}
    \mathcal{L} = - \frac{1}{B} \sum_{i=1}^{B} \sum_{x \in X} \log \frac{1}{1 + e^{-t_i(x) \left( s \inner{f(q_i)}{f(x)} + b \right)}} ,
\end{equation}

where $X$ is the set of $2B$ candidate documents in the batch, $t_i(x) = +1$ if $x = x_i^+$ and $-1$ otherwise, and $s$ and $b$ are a learnable scale and bias.
Each pair is fit to its own target, so the loss constrains the absolute score of every candidate rather than its rank.

% CoSENT \citep{huang2024cosent} consumes pairs $(q_i, x_i)$ with relevance labels $r_i$ and penalizes every two pairs whose similarities are ordered against their labels:
CoSENT \citep{huang2024cosent} consumes pairs $(q_i, x_i)$ with ordinal relevance labels $r_i$ and penalizes every two pairs whose similarities violate ordering:

\begin{equation}
    \label{eq:trad-cosent}
    \mathcal{L} = \log \Big( 1 + \sum_{r_i > r_j} \exp\big( s \big( \inner{f(q_j)}{f(x_j)} - \inner{f(q_i)}{f(x_i)} \big) \big) \Big) ,
\end{equation}

where the sum runs over all ordered pairs of pairs in the batch with $r_i > r_j$.
The labels enter only through their ordering, so with binary labels $r_i \in \{0, 1\}$ the loss requires every positive pair to outscore every negative pair by a margin set by $s$.


\section{NV-Retriever Implementation Details}
\label{app:nv-sweep}

\darkgreen{We mine the train split only, with \texttt{e5-mistral-7b-instruct} for text and \texttt{clip-ViT-L-14} for images, retrieving $1000$ candidates per text query and $100$ per image query before the positive-aware filter; if no candidate survives, we retain the lowest-scoring retrieved non-positive candidate.}
We grid-search the filter's relative margin $m \in \{0.025, 0.05, 0.1, 0.2\}$ against a survivor offset $s \in \{0, 10\}$, taking the $(s+1)$-th survivor rather than the first.
\darkgreen{If fewer than $s+1$ candidates survive, we take the last survivor.}


\section{List of Prompts}
\label{app:prompts}

Placeholders in braces are substituted at generation time.
Attribute counts appear resolved to their configured value of five per bucket.

\subsection{Product Attribute Extraction}
\label{prompt:text-features}

Applied to each (positive, hard negative) product pair drawn from ESCI, with both product
records serialized into the corresponding placeholders.

\begin{lstlisting}[style=promptstyle]
### Product A:

{positive_product}

### Product B:

{substitute_irrelevant_product}

List up to 5 features common to both products ("common_features"), up to 5 features unique to product A ("unique_features_a"), up to 5 features unique to product B ("unique_features_b"), and up to 5 features that do not apply to either product ("neither_features")). When generating common_features, ensure they are present in both products in some way. For example, if product A is "low fat" and product B in "low sugar", then "low calories" could be in common_features. If there are fewer than 5 common features, generate as many as possible. Do not assume anything about products not written in the description. When generating neither_features, ensure they are opposite or mutually exclusive with features in one or both of the products. For example, if product A is "red" and product B is "orange", then "blue" could be in neither_features. Must not generate negated features, such as "not red", "no wifi", "without bluetooth", or "doesn't have x".Dont use "not" or "without". Use your imagination and create diverse neither_features. All features should be objective and no more than 5 words. Return ONLY JSON: {'common_features': ['feature1', 'feature2', ...], 'unique_features_a': ['feature1', 'feature2', ...], 'unique_features_b': ['feature1', 'feature2', ...], 'neither_features': ['feature1', 'feature2', ...]}
\end{lstlisting}

\subsection{Visual Attribute Extraction}
\label{prompt:image-features}

Applied to each (positive, hard negative) garment pair drawn from DeepFashion In-Shop.
Both photographs are supplied to the model alongside the prompt, in the order named.

\begin{lstlisting}[style=promptstyle]
### Product A:

The first image is Product A. Metadata: {positive_product}

### Product B:

The second image is Product B. Metadata: {hard_negative_product}

List up to 5 visual features common to both products ("common_features"), up to 5 visual features unique to product A ("unique_features_a"), up to 5 visual features unique to product B ("unique_features_b"), and up to 5 visual features that do not apply to either product ("neither_features").

When generating common_features, ensure they are visible in both products in some way. If there are fewer than 5 common features, generate as many as possible. Do not assume anything not visible in the images or stated in metadata.

When generating neither_features, ensure they are opposite or mutually exclusive with visual features in one or both products. For example, if product A is red and product B is orange, then "blue" could be in neither_features. Must not generate negated features, such as "not red", "no sleeves", "without pattern", or "doesn't have x". Dont use "not" or "without". Use your imagination and create diverse neither_features. All features should be objective and no more than 5 words.

Return ONLY JSON: {"common_features": ["feature1", "feature2"], "unique_features_a": ["feature1", "feature2"], "unique_features_b": ["feature1", "feature2"], "neither_features": ["feature1", "feature2"]}
\end{lstlisting}

\subsection{Query Rephrasing}
\label{prompt:rephrase}

\darkgreen{Applied to the templated query with the positive product's text and five modality-specific human-written example queries to produce the rewritten queries used in the Main dataset.}

\begin{lstlisting}[style=promptstyle,basicstyle=\ttfamily\scriptsize\color{darkgreen}]
Original query: '{query}'

Context: The user will be happy with the following product:

{product_description}

Task: Rephrase the query so it keeps the same meaning but sounds more natural and uses synonyms.
Rephrase it similar to the style of these example queries, but avoid being exactly like them:

1. {example_query_1}
2. {example_query_2}
3. {example_query_3}
4. {example_query_4}
5. {example_query_5}

Rules:
- Preserve intent exactly (positive stays positive, negative stays negative).
- You MUST rephrase at least half of the comma-separated attributes.
- Do NOT force rewording if it sounds unnatural.
- Do NOT change the product/item name unless a very natural alternative exists.
- Do NOT change measurement units or numeric values.
- Keep the query fluent and realistic.

Return ONLY a JSON object: {"rephrased_query": "..."}
\end{lstlisting}



\begingroup

\section{Annotator Instructions}
\label{app:annotator-instructions}

Welcome to the Columbia Advanced Search Study.
In this study, you will be helping us generate complex search phrasings for product and image search.

\begin{itemize}
    \item Do not include any personally identifiable information in your responses
    \item Do not use AI (ChatGPT, Claude, Gemini, etc.) for this task
\end{itemize}

If you feel discomfort in this task, you can discontinue at any time.

\par\medskip\noindent\rule{\linewidth}{0.4pt}\par\medskip

Below is an example product description, followed by the questions for it.

Imagine you are shopping online for an \textbf{ipad 6 gen case.}
Read the product description and answer the following questions about it.

[Product Description 1]

[Product Description 2]

\begin{enumerate}
    \item List n\_1 attributes of \textbf{Product 1} described in the product description that you like about it.
    These can be quotes from the product description, up to 8 words long.
    The attributes must be objective (red \(\checkmark\), cute \(\oslash\)).
    Separate your attributes with commas.
    If your example contains an image, make sure the features you are describing are visible in the image.

    Example response: Anti-fingerprint coating, easy access to all buttons, fits iPad Air 2

    \item List n\_2 attributes that \textbf{Product 1 does not have} that you \textbf{do not} want it to have.
    These can be attributes of Product 2, but do not have to be.
    Separate your attributes with commas.

    Example response: Graffiti pattern, Anti-slip grooves, hands-free viewing

    \item Imagine you are searching Amazon for a product that has all of the attributes listed in question (1) and none of the attributes listed in question (2).
    Write a search as if you were searching Amazon.
    Make sure to include the item itself (e.g. ipad 6 gen case).
    It is \textbf{okay} to include exact quotes from the product

    Example response: 6 gen ipad case with anti-fingerprint coating, easy access to all buttons fitting iPad Air 2.
    No graffiti pattern, anti-slip grooves, or hands free viewing.

    \item Rewrite the query in (3) to avoid exact quotes from either product description, if possible.

    Example response: 6 gen ipad case with non-glossy coating and accessible buttons for iPad Air 2.
    Does not have any styling, grooves, or stand.
\end{enumerate}

\endgroup

\end{document}
